\documentclass[a4paper]{arrowhead}

\usepackage[yyyymmdd]{datetime}
\usepackage{etoolbox}
\usepackage[utf8]{inputenc}
\usepackage[T1]{fontenc}
\usepackage{multirow}
\usepackage{longtable}
\usepackage{tikz}
\usepackage{adjustbox}
\AtBeginEnvironment{longtable}{\small}
\BeforeBeginEnvironment{tikzpicture}{\begin{adjustbox}{max width=\textwidth}}
\AfterEndEnvironment{tikzpicture}{\end{adjustbox}}
\usetikzlibrary{positioning,shapes.multipart,arrows.meta}

\renewcommand{\dateseparator}{-}
\setlength{\emergencystretch}{3em}
\lstset{language=HTTP,breakatwhitespace=false,columns=fullflexible}

\newcommand{\fparam}[1]{\textit{\textcolor{ArrowheadBlue}{#1}}}

%% Special references
\newcommand{\mref}[2][]{{\textcolor{ArrowheadPurple}{\hyperref[sec:model:#2]{\ifstrempty{#1}{#2}{#1}}}}}
\newcommand{\oref}[1]{{\textcolor{ArrowheadBlue}{\hyperref[sec:operations:#1]{#1}}}}
\newcommand{\pdef}[1]{{\textcolor{ArrowheadGrey}{#1 \label{sec:model:primitives:#1} \label{sec:model:primitives:#1s}}}}
\newcommand{\pref}[1]{{\textcolor{ArrowheadGrey}{\hyperref[sec:model:primitives:#1]{#1}}}}

% as template_idd.tex; the input and output arguments are passed already formatted (\mref or \pref)
\newrobustcmd\osubsection[5]{
  \addtocounter{subsection}{1}
  \addcontentsline{toc}{subsection}{\protect\numberline{\thesubsection}operation \textcolor{ArrowheadBlue}{#1}}
  \subsection*{
    \thesubsection\quad
    #2 \textcolor{ArrowheadPurple}{#3} \\
    \small
    \hspace*{0.075\textwidth}\begin{minipage}{0.1\textwidth}
      \vspace*{1mm}
      Operation: \\
      \notblank{#4}{Input: \\}{}
      \notblank{#5}{Output: \\}{}
    \end{minipage}%
    \begin{minipage}{0.825\textwidth}
      \vspace*{1mm}
      \textcolor{ArrowheadBlue}{#1} \\
      \notblank{#4}{#4 \\}{}
      \notblank{#5}{#5 \\}{}
    \end{minipage}
  }
  \label{sec:operations:#1}
}
\newrobustcmd\msubsection[2]{
  \addtocounter{subsection}{1}
  \addcontentsline{toc}{subsection}{\protect\numberline{\thesubsection}#1 \textcolor{ArrowheadPurple}{#2}}
  \subsection*{\thesubsection\quad#1 \textcolor{ArrowheadPurple}{#2}}
  \label{sec:model:#2}
}
%%

\begin{document}
\sloppy

%% Arrowhead Document Properties
\ArrowheadTitle{deviceDiscovery MQTT/TLS/JSON}
\ArrowheadServiceID{deviceDiscovery}
\ArrowheadType{Interface Design Description}
\ArrowheadTypeShort{IDD}
\ArrowheadVersion{5.2.2}
\ArrowheadDate{2026-10-22}
\ArrowheadAuthor{Jerker Delsing}
\ArrowheadStatus{RELEASE CANDIDATE}
\ArrowheadContact{jerker.delsing@ltu.se}
\ArrowheadFooter{\href{www.arrowhead.eu}{www.arrowhead.eu}}
\ArrowheadSetup
%%

%% Front Page
\begin{center}
  \vspace*{1cm}
  \huge{\arrowtitle}

  \vspace*{0.2cm}
  \LARGE{\arrowtype}
  \vspace*{1cm}

  \Large{Service ID: \textit{"\arrowid"}}

  \vspace*{\fill}

  % Front Page Image

  \vspace*{1cm}
  \vspace*{\fill}

  % Front Page Abstract
  \begin{abstract}
    This document describes the MQTT protocol with TLS payload security and JSON payload encoding variant of the \texttt{deviceDiscovery} service of the ServiceRegistry system.
  \end{abstract}

  \vspace*{1cm}
\end{center}

\newpage
%%

%% Table of Contents
\tableofcontents
\newpage
%%

\section{Overview}
\label{sec:overview}

This document describes the \texttt{deviceDiscovery} service interface of the ServiceRegistry system over MQTT, which provides registration and lookup of the devices of the local cloud: the physical or virtual hosts, with their network addresses and metadata, on which the systems run.

It's implemented using protocol, encoding as stated in below table:
\begin{table}[ht!]
  \centering
  \begin{tabular}{|l|l|l|}
    \hline
    \rowcolor{gray!33} Profile type & Type & Version \\ \hline
    Transfer protocol & MQTT & 3.1.1 \cite{mqtt311} \\ \hline
    Data encryption & TLS & 1.3 \cite{rfc8446}, 1.2 \cite{rfc5246} \\ \hline
    Encoding & JSON & RFC 8259 \cite{rfc8259} \\ \hline
    Compression & N/A & - \\ \hline
    Semantics & N/A & - \\ \hline
    Ontology & N/A & - \\ \hline
  \end{tabular}
  \caption{Communication and semantics details used for the \texttt{deviceDiscovery} service interface}
  \label{tab:comunication_semantics_profile}
\end{table}

This document provides the Interface Design Description IDD to the \textit{deviceDiscovery -- Service Description} document \cite{sr-deviceDiscovery-sd}.
The operations, their semantics and the data models are the same as in the HTTP/TLS/JSON IDD \cite{sr-deviceDiscovery-http-idd}; this document specifies the MQTT binding.

The MQTT interface is offered only when the system runs with \texttt{mqtt.api.enabled=true}; it is then registered with the interface template \texttt{generic\_\allowbreak{}mqtts} (\texttt{generic\_\allowbreak{}mqtt} without TLS), the broker address and port (\texttt{mqtt.\allowbreak{}broker.\allowbreak{}address}, \texttt{mqtt.\allowbreak{}broker.\allowbreak{}port}) and the base topic \texttt{arrowhead/\allowbreak{}serviceregistry/\allowbreak{}device-\allowbreak{}discovery/\allowbreak{}}. In the Arrowhead-test-cloud-java deployment the MQTT interfaces of the core systems are disabled.

\newpage

\section{Service Operations}
\label{sec:operations}

The provider subscribes to one topic per operation, the base topic followed by the operation name:

\noindent\begin{longtable}{| p{5cm} | p{10.4cm} |} \hline
\rowcolor{gray!33} Operation & Request topic \\ \hline \endhead
\texttt{register} & \texttt{arrowhead/\allowbreak{}serviceregistry/\allowbreak{}device-\allowbreak{}discovery/\allowbreak{}register} \\ \hline
\texttt{lookup} & \texttt{arrowhead/\allowbreak{}serviceregistry/\allowbreak{}device-\allowbreak{}discovery/\allowbreak{}lookup} \\ \hline
\texttt{revoke} & \texttt{arrowhead/\allowbreak{}serviceregistry/\allowbreak{}device-\allowbreak{}discovery/\allowbreak{}revoke} \\ \hline
\end{longtable}

A consumer publishes a request envelope (\mref{MqttRequestTemplate}) to the topic and receives the answer (\mref{MqttResponseTemplate}) on its own \texttt{responseTopic}. The \texttt{authentication} field carries the requester's identity according to the provider's authentication policy: \texttt{IDENTITY-TOKEN//<token>} (outsourced or internal policy), \texttt{SYSTEM//<name>} (declared policy) or the requester's X.509 certificate (certificate policy). The \texttt{status} of the response uses the HTTP status code numbers (200, 201, 204, 400, 401, 403, 404, 408, 423, 500, 503); on error the payload is an \mref{ErrorMessage}.

\begin{lstlisting}[caption={Request envelope.}]
{
  "traceId": "a1b2c3",
  "authentication": "IDENTITY-TOKEN//<token>",
  "responseTopic": "arrowhead/response/TemperatureDisplay/a1b2c3",
  "qosRequirement": 1,
  "params": {},
  "payload": {}
}
\end{lstlisting}

\begin{lstlisting}[caption={Response envelope.}]
{
  "status": 200,
  "traceId": "a1b2c3",
  "receiver": "TemperatureDisplay",
  "payload": {}
}
\end{lstlisting}

\osubsection{register}{PUBLISH}{\texttt{arrowhead/\allowbreak{}serviceregistry/\allowbreak{}device-\allowbreak{}discovery/\allowbreak{}register}}{\mref{DeviceRequest}}{\mref{DeviceResponse}}

Registers a device with its addresses and metadata. If a device with the same name and the same data exists, it is returned unchanged (200); a new device is created (201). A device with the same name and different data is refused.

\begin{lstlisting}[caption={Payload of an \oref{register} request.}]
{
  "name": "Raspberry1",
  "metadata": {
    "location": "hall-1"
  },
  "addresses": [
    "string"
  ]
}
\end{lstlisting}

\osubsection{lookup}{PUBLISH}{\texttt{arrowhead/\allowbreak{}serviceregistry/\allowbreak{}device-\allowbreak{}discovery/\allowbreak{}lookup}}{\mref{DeviceLookupRequest}}{\mref{DeviceListResponse}}

Returns the devices matching all given filters; without filters all devices.

\begin{lstlisting}[caption={Payload of an \oref{lookup} request.}]
{
  "deviceNames": [
    "Raspberry1"
  ],
  "addresses": [
    "string"
  ],
  "addressType": "IPV4",
  "metadataRequirementList": [
    {
      "unit": "celsius"
    }
  ]
}
\end{lstlisting}

\osubsection{revoke}{PUBLISH}{\texttt{arrowhead/\allowbreak{}serviceregistry/\allowbreak{}device-\allowbreak{}discovery/\allowbreak{}revoke}}{\pref{String}}{}

Removes the named device. Returns 200 if it was removed, 204 if it did not exist, and 423 Locked if systems are still assigned to it.

\begin{lstlisting}[caption={Payload of an \oref{revoke} request.}]
"TemperatureSensor"
\end{lstlisting}

\newpage

\section{Data Models}
\label{sec:model}

Here, all data objects that can be part of the service calls associated with this service are listed in alphabetic order.
Each \textit{struct} denotes a JSON \pref{Object} with the named fields; fields that are not needed may be omitted.

\msubsection{struct}{MqttRequestTemplate}

The request envelope.

\noindent\begin{tabularx}{\textwidth}{| p{4.25cm} | p{3.5cm} | X |} \hline
\rowcolor{gray!33} Object Field & Value Type & Description \\ \hline
"traceId" & \pref{String} & Identifier chosen by the requester, returned in the response. \\ \hline
"authentication" & \pref{String} & Identity of the requester (see Section \ref{sec:operations}). \\ \hline
"responseTopic" & \pref{String} & Topic on which the response is published. \\ \hline
"qosRequirement" & \pref{Number} & MQTT QoS of the response (0, 1 or 2). \\ \hline
"params" & \pref{Object}$<$\pref{String}$>$ & Optional parameters (the query parameters of the HTTP binding). \\ \hline
"payload" & \pref{Object} & The operation input. \\ \hline
\end{tabularx}

\msubsection{struct}{MqttResponseTemplate}

The response envelope.

\noindent\begin{tabularx}{\textwidth}{| p{4.25cm} | p{3.5cm} | X |} \hline
\rowcolor{gray!33} Object Field & Value Type & Description \\ \hline
"status" & \pref{Number} & Status code. \\ \hline
"traceId" & \pref{String} & The trace identifier of the request. \\ \hline
"receiver" & \pref{String} & Name of the requester. \\ \hline
"payload" & \pref{Object} & The operation output, or an \mref{ErrorMessage}. \\ \hline
\end{tabularx}

\msubsection{struct}{ErrorMessage}




\noindent\begin{tabularx}{\textwidth}{| >{\raggedright\arraybackslash}p{3.9cm} | >{\raggedright\arraybackslash\small}p{5.3cm} | X |} \hline
\rowcolor{gray!33} Object Field & Value Type & Description \\ \hline
"error\allowbreak{}Message" & \pref{String} &  \\ \hline
"error\allowbreak{}Code" & \pref{Number} &  \\ \hline
"exception\allowbreak{}Type" & \pref{ExceptionType} &  \\ \hline
"origin" & \pref{String} &  \\ \hline
\end{tabularx}



\msubsection{struct}{Address}

A network address with its type.


\noindent\begin{tabularx}{\textwidth}{| >{\raggedright\arraybackslash}p{3.9cm} | >{\raggedright\arraybackslash\small}p{5.3cm} | X |} \hline
\rowcolor{gray!33} Object Field & Value Type & Description \\ \hline
"type" & \pref{String} & Address type: \texttt{HOSTNAME}, \texttt{IPV4}, \texttt{IPV6} or \texttt{MAC}. \\ \hline
"address" & \pref{String} & One network address (host name, IPv4, IPv6 or MAC address). \\ \hline
\end{tabularx}



\msubsection{struct}{DeviceListResponse}

List of devices.


\noindent\begin{tabularx}{\textwidth}{| >{\raggedright\arraybackslash}p{3.9cm} | >{\raggedright\arraybackslash\small}p{5.3cm} | X |} \hline
\rowcolor{gray!33} Object Field & Value Type & Description \\ \hline
"entries" & \pref{Array}$<$\allowbreak{}\mref{DeviceResponse}$>$ & Matching records of the current page. \\ \hline
"count" & \pref{Number} & Total number of matching records (all pages). \\ \hline
\end{tabularx}



\msubsection{struct}{DeviceLookupRequest}

Filter for devices; all fields optional, combined with AND.


\noindent\begin{tabularx}{\textwidth}{| >{\raggedright\arraybackslash}p{3.9cm} | >{\raggedright\arraybackslash\small}p{5.3cm} | X |} \hline
\rowcolor{gray!33} Object Field & Value Type & Description \\ \hline
"device\allowbreak{}Names" & \pref{Array}$<$\allowbreak{}\pref{String}$>$ & Device names to match. \\ \hline
"addresses" & \pref{Array}$<$\allowbreak{}\pref{String}$>$ & Network addresses of the entity. \\ \hline
"address\allowbreak{}Type" & \pref{String} & Address type filter: \texttt{HOSTNAME}, \texttt{IPV4}, \texttt{IPV6} or \texttt{MAC}. \\ \hline
"metadata\allowbreak{}Requirement\allowbreak{}List" & \pref{Array}$<$\allowbreak{}\mref{MetadataRequirement}$>$ & Metadata requirements; each map is one alternative. \\ \hline
\end{tabularx}



\msubsection{struct}{DeviceRequest}

A device and its addresses.


\noindent\begin{tabularx}{\textwidth}{| >{\raggedright\arraybackslash}p{3.9cm} | >{\raggedright\arraybackslash\small}p{5.3cm} | X |} \hline
\rowcolor{gray!33} Object Field & Value Type & Description \\ \hline
"name" & \pref{String} & Device name, PascalCase. \\ \hline
"metadata" & \pref{Object} & Free-form metadata (JSON object). \\ \hline
"addresses" & \pref{Array}$<$\allowbreak{}\pref{String}$>$ & Network addresses of the entity. \\ \hline
\end{tabularx}



\msubsection{struct}{DeviceResponse}

A registered device.


\noindent\begin{tabularx}{\textwidth}{| >{\raggedright\arraybackslash}p{3.9cm} | >{\raggedright\arraybackslash\small}p{5.3cm} | X |} \hline
\rowcolor{gray!33} Object Field & Value Type & Description \\ \hline
"name" & \pref{String} & Device name. \\ \hline
"metadata" & \pref{Object} & Free-form metadata (JSON object). \\ \hline
"addresses" & \pref{Array}$<$\allowbreak{}\mref{Address}$>$ & Network addresses with their types. \\ \hline
"created\allowbreak{}At" & \pref{DateTime} & UTC time of creation. \\ \hline
"updated\allowbreak{}At" & \pref{DateTime} & UTC time of the last change. \\ \hline
\end{tabularx}



\msubsection{struct}{MetadataRequirement}

A JSON object of metadata requirements: each key is a metadata path (dot notation), each value the required value or an operator object such as \texttt{\{"op": "LESS\_\allowbreak{}THAN",\allowbreak{} "value": 10\}}.



\subsection{Primitives}
\label{sec:model:primitives}

As all messages are encoded using the JSON format \cite{rfc8259}, the following primitive constructs, part of that standard, become available.

\noindent\begin{tabularx}{\textwidth}{| p{4.6cm} | X |} \hline
\rowcolor{gray!33} Type & Description \\ \hline
\pdef{Array} & A JSON Array \cite{rfc8259}; \pref{Array}$<$A$>$ has elements of type A. \\ \hline
\pdef{DateTime} & JSON String with an ISO 8601 UTC timestamp, e.g. \texttt{2026-10-22T12:00:00Z}. \\ \hline
\pdef{ExceptionType} & One of \texttt{ARROWHEAD}, \texttt{INVALID\_\allowbreak{}PARAMETER}, \texttt{AUTH}, \texttt{FORBIDDEN}, \texttt{DATA\_\allowbreak{}NOT\_\allowbreak{}FOUND}, \texttt{TIMEOUT}, \texttt{LOCKED}, \texttt{INTERNAL\_\allowbreak{}SERVER\_\allowbreak{}ERROR}, \texttt{EXTERNAL\_\allowbreak{}SERVER\_\allowbreak{}ERROR} (JSON String). \\ \hline
\pdef{Number} & JSON Number, an integer. \\ \hline
\pdef{Object} & JSON Object; \pref{Object}$<$A$>$ has string keys and values of type A. \\ \hline
\pdef{String} & JSON String. \\ \hline
\end{tabularx}

\newpage

\bibliographystyle{IEEEtran}
\bibliography{bibliography}

\newpage

\section{Revision History}
\subsection{Amendments}

\noindent\begin{tabularx}{\textwidth}{| p{1cm} | p{3cm} | p{2cm} | X | p{4cm} |} \hline
\rowcolor{gray!33} No. & Date & Version & Subject of Amendments & Author \\ \hline

1 & 2026-10-22 & \arrowversion & Document created from the Arrowhead LaTeX templates for the 5.2.2 Java Spring implementation, including the local extensions of this workspace. & Jerker Delsing \\ \hline
\end{tabularx}

\subsection{Quality Assurance}

\noindent\begin{tabularx}{\textwidth}{| p{1cm} | p{3cm} | p{2cm} | X |} \hline
\rowcolor{gray!33} No. & Date & Version & Approved by \\ \hline

1 & & \arrowversion  &  \\ \hline

\end{tabularx}

\end{document}
